\documentclass[10pt,a4paper]{amsart}
\usepackage[margin=19mm]{geometry}
\usepackage{amsmath,amssymb,mathtools}
\usepackage[scr=rsfs,cal=euler]{mathalfa}
\usepackage{xfrac}
\usepackage[unicode,hidelinks,bookmarks=false]{hyperref}
\newcommand{\ie}{i.e.\ }
\newcommand{\cf}{cf.\ }
\newcommand{\resp}{respectively\ }
\newcommand{\Subsection}{\subsection}
\providecommand{\nobreakdash}{\nobreak\hbox{-}}
\setlength{\emergencystretch}{2em}
\multlinegap0pt
\usepackage[shortlabels]{enumitem}
\usepackage{dsfont}
\usepackage{amssymb}
\usepackage{xcolor}
\usepackage{tikz}
\newtheorem{theorem}{Theorem}
\newtheorem{claim}{Claim}
\newcommand{\cF}{\mathcal{F}}
\title{Upper bounds on the running time of bootstrap percolation}
\author{Weichan Liu, Xiangxiang Nie, Sim\'on Piga, Bjarne Sch\"{u}lke}
\date{Source: arXiv:2604.22630v2 (CC BY 4.0)}
\begin{document}
\maketitle

\noindent\emph{Source and licence.} Upper bounds on the running time of bootstrap percolation. Version arXiv:2604.22630v2 (CC BY 4.0), distributed by arXiv under the Creative Commons Attribution 4.0 International licence, http://creativecommons.org/licenses/by/4.0/, as recorded in the arXiv licence metadata for this exact version and shown on the abstract page https://arxiv.org/abs/2604.22630v2. Full licence notice: \texttt{licenses/arxiv-2604.22630-CC-BY-4.0.txt}.\par\smallskip

\noindent\textit{Editorial scope.} Counted unit: the article's theorem thm:main (source0.tex lines 716--723). The article's complete original proof of it is delivered with it (source0.tex lines 782--880, one \texttt{proof} environment of 99 lines, which the article places away from the statement). No mathematics is written, repaired or re-derived: the statement and the proof are the article's own lines, in the article's own notation, and no retained line is substituted or re-flowed.  No further source-local context is needed: every cross-reference of every retained line resolves inside the two delivered spans.  Nothing was omitted from any retained span, so the package carries no omission record.  No retained line contains a citation, so the wrapper carries no bibliography and no bibliography entry.  No retained line contains a figure environment, an image-inclusion command or drawing code, so no image is delivered and the package declares no asset dependency.  Editorial formatting only, in addition: an AMS \texttt{amsart} preamble that restates the article's own declarations and macros used by the retained lines, namely the environments \texttt{theorem}, \texttt{claim} on separate counters, together with the standard packages those lines need.  The retained environments print in this document as Theorem 1, Claim 1 in order of appearance, whereas the article prints the same items with the numbering of its own full document.  The complete native source of the article as distributed in the arXiv e-print for this exact version is bundled unchanged as \texttt{sources/199112/source0.tex}.
    \begin{theorem}\label{thm:main}
        For every integer~$k\geq2$, every~$k$-graph~$F$, and every~$e\in E(F)$ we have 
        \begin{equation*}
         M_F(n)
         \leq
         \big(\pi(F-e)+o(1)\big)\binom{n}{k}\,.
        \end{equation*}
    \end{theorem}

    \begin{proof}[Proof of Theorem~\ref{thm:main}]
        Given~$\varepsilon>0$, an integer~$k\geq2$, and a~$k$-graph~$F$, choose~$\delta>0$ and~$n\in\mathds{N}$ such that $$\varepsilon,k,v(F)\gg\delta\gg n^{-1}\,.$$
        Let~$e\in E(F)$ and set~$F^-=F-e$.
        We need to show that~$M_F(n)
         \leq
         \big(\pi(F^-)+\varepsilon\big)\binom{n}{k}$.
        One can observe that the statement holds if~$F$ has at most two edges, so assume that~$F$ has at least three edges.
        Let~$(H_i)_{i\geq0}$ be an~$F$-process starting with~$H=(V,E)$, where~$\vert V\vert=n$, that satisfies~$M_F(H)= M_F(n)=:\tau$.
        For each~$i\in[\tau]$ choose~$e_i\in E(H_i)\setminus E(H_{i-1})$ arbitrarily.
        Let~$G$ be the graph on~$V$ whose edge set is~$\{e_i:\:i\in[\tau]\}$.
        Assume that~$\tau\geq(\pi(F^-)+\varepsilon)\binom{n}{k}$.
        We aim to derive a contradiction by establishing upper and lower bounds for the number of copies of~$F^-$ in~$G$.
        Denote the set of copies of~$F^-$ in~$G$ by~$\cF$.
        Supersaturation and the hierarchy of constants imply that
        \begin{align}\label{eq:Flower}
            \vert\cF\vert\geq\delta n^{v(F^-)}=\delta n^{v(F)}\,.
        \end{align}
        For every~$D\in\cF$ define~$i(D)=\max\{i\in[\tau]:\:e_i\in E(D)\}$ and~$j(D)=\max\{j\in[\tau]:\:e_j\in E(D)\setminus e_{i(D)}\}$.
        Furthermore, for every~$D\in\cF$ fix some~$e_D\in V^{(k)}$ such that~$D \cup e_D$ is a copy of~$F$ (if there are several choices, pick one arbitrarily).
        For~$e\in E(H_{\tau})$ let~$\mathfrak{s}(e)$ denote the step in which~$e$ is added, i.e.,~$\mathfrak{s}(e)=\min\{s\in[\tau]:\:e\in E(H_s)\setminus E(H_{s-1})\}$ if~$e\notin E(H_0)$ and~$\mathfrak{s}(e)=0$ if~$e\in E(H_0)$.
        Depending on how~$\mathfrak{s}(e_D)$ relates to~$j(D)$ and~$i(D)$, we assign to every~$D \in \cF$ one of the following four types.
        \begin{enumerate}[leftmargin=1.95cm]
         %\addtolength{\leftmargin}{em}
            \item[Type~$1$:]\label{it:type1}
            $\mathfrak{s}(e_D)\leq j(D)$,
            \item[Type~$2$:]\label{it:type2}
            $j(D)< \mathfrak{s}(e_D)< i(D)$,
            \item[Type~$3$:]\label{it:type3}
            $i(D) < \mathfrak{s}(e_D)$,
            \item[Type~$4$:]\label{it:type4}
            $\mathfrak{s}(e_D)=i(D)$.
        \end{enumerate}

    Let~$\cF_{\alpha}$ be the family of copies~$D\in \cF$ of Type~$\alpha$ for~$\alpha\in[4]$.
    Further, for~$i\in[\tau]$ let~$\cF_{\alpha}^i=\{D\in\cF_{\alpha}:\:i(D)=i\}$.
    Next we find upper bounds on~$\vert\cF_{\alpha}\vert$ for every~$\alpha\in[4]$.

    \begin{claim}\label{cl:F1}
        $\vert \cF_1\vert\leq n^{v(F)-1}$.
    \end{claim}

    \begin{proof}
        First observe that for~$D \in \cF_1$ we must have~$e_{j(D)} = e_{i(D)-1}$.
        Indeed, from the definition of Type~$1$ and~$G$ it follows that in~$H_{j(D)}\cup e_{i(D)}$ there is a copy of~$F$ containing~$e_{i(D)}$.
        Hence,~$e_{i(D)}\in E(H_{j(D)+1})$.
        Due to the definition of~$G$, this means that~$i(D)=j(D)+1$.
        
        Now fix~$j\in[\tau]$ with~$j\geq2$ and let~$v \in e_{j-1} \setminus e_j$. 
        By the above, every~$D \in \mathcal{F}_1^j$ must contain~$v$. 
        In particular, every~$D\in\cF_1^j$ contains the~$k+1$ vertices in the set~$e_j\cup v$.
        Since no~$D\in\cF$ satisfies~$i(D)=1$, it thus follows that for every~$j\in[\tau]$ we have~$
        \vert\cF_1^j\vert\leq n^{v(F)-(k+1)}$. 
        Summing over all~$j\in[\tau]$, we obtain
        \begin{align*}
            \vert\mathcal{F}_1\vert=\sum_{j\in[\tau]} \vert\cF_1^j\vert\leq \binom{n}{k} \cdot n^{v(F)-k-1}= n^{v(F)-1}\,,
        \end{align*}
        as desired.
    \end{proof}
    
    To show the remaining upper bounds, we first observe the following for~$\alpha\in\{2,3,4\}$.
    Since for every~$D\in\cF_{\alpha}$, it holds that~$|e_D \cup e_{i(D)}|\geq k+1$ and~$e_D \cup e_{i(D)}\subseteq V(D)$, we have~$\vert\{D\in\cF_{\alpha}^i:\:e_D=e\}\vert\leq n^{v(F)-\vert e_D\cup e_{i(D)}\vert}\leq n^{v(F)-k-1}$ for every~$e\in V^{(k)}$.
    Further, note the trivial fact that if for some~$e\in V^{(k)}$ there is no~$D\in\cF_{\alpha}^i$ with~$e_D=e$, then~$\vert\{D\in\cF_{\alpha}^i:\:e_D=e\}\vert=0$.
    Thus, we infer
    \begin{align}\label{eq:generalupper}
        |\cF_{\alpha}|&=\sum_{i\in[\tau]}\vert\cF_{\alpha}^i\vert=\sum_{i \in [\tau]} \sum_{e\in V^{(k)}}\vert\{D\in\cF_{\alpha}^i:\:e_D=e\}\vert\nonumber\\
        &\leq\sum_{i \in [\tau]}\big\vert\{e\in V^{(k)}:\:e=e_D\text{ for some }D\in\cF_{\alpha}^i\}\big\vert \cdot n^{v(F) - (k+1)}\,.
    \end{align}

    Next we show that every~$e\in V^{(k)}$ can only be counted once in the above sum.
    \begin{claim}\label{cl:disjointsets}
        For every~$\alpha\in\{2,3,4\}$ and distinct~$i,j\in[\tau]$ we have $$\{e\in V^{(k)}:\:e=e_D\text{ for some }D\in\cF_{\alpha}^i\}\cap\{e\in V^{(k)}:\:e=e_D\text{ for some }D\in\cF_{\alpha}^j\}=\emptyset\,.$$
    \end{claim}

    \begin{proof}
        Observe that for every~$D\in\cF_2$ we must have~$\mathfrak{s}(e_D)=i(D)-1$.
        Indeed, by the definition of Type~$2$ and~$G$, we know that in~$H_{\mathfrak{s}(e_D)}\cup e_{i(D)}$ there is a copy of~$F$ containing~$e_{i(D)}$.
        Keeping in mind that~$\mathfrak{s}(e_D)<i(D)$ and the definition of the~$F$-process, we conclude~$i(D)=\mathfrak{s}(e_D)+1$.
        Therefore, once we fix~$i\in[\tau]$, we know that for every~$D\in\cF_2^i$ we must have~$\mathfrak{s}(e_D)=i-1$.
        This in turn entails that the sets~$\{e\in V^{(k)}:\:e=e_D\text{ for some }D\in\cF_2^i\}$ are disjoint for distinct~$i\in[\tau]$.
        
        Note that for every~$D\in\cF_3$ we know that~$H_{i(D)}\cup e_D$ contains a copy of~$F$ with~$e_D$ being an edge of this copy, and that~$e_D\notin E(H_{i(D)})$ by the definition of Type~$3$.
        Hence, the definition of the~$F$-process implies that~$e(D)$ is added in the next step, i.e.,~$\mathfrak{s}(e_D)=i(D)+1$.
        Therefore, once we fix~$i\in[\tau]$, we know that for every~$D\in\cF_3^i$ we must have~$\mathfrak{s}(e_D)=i+1$.
        This in turn entails that the sets~$\{e\in V^{(k)}:\:e=e_D\text{ for some }D\in\cF_3^i\}$ are disjoint for distinct~$i\in[\tau]$.
        
        For~$D\in\cF_4$, we have~$\mathfrak{s}(e_D)=i(D)$ by definition of Type~$4$.
        Therefore, once we fix~$i\in[\tau]$, we know that every~$D\in\cF_4^i$ must satisfy~$\mathfrak{s}(e_D)=i$ and the claim follows.
    \end{proof}

    Claim~\ref{cl:disjointsets} reveals that for~$\alpha\in\{2,3,4\}$ we have $$\sum_{i \in [\tau]}\big\vert\{e\in V^{(k)}:\:e=e_D\text{ for some }D\in\cF_{\alpha}^i\}\big\vert\leq\binom{n}{k}\,.$$
    Combining this with~\eqref{eq:generalupper} yields
    \begin{align}\label{eq:Falpha}
        \vert\cF_{\alpha}\vert\leq n^{v(F)-(k+1)}\cdot\binom{n}{k}\leq n^{v(F)-1}\,.
    \end{align}

    Comparing the bounds in~\eqref{eq:Falpha} and the bound in Claim~\ref{cl:F1} with~\eqref{eq:Flower} entails
    $$\delta n^{v(F)}\leq\vert\cF\vert\leq\sum_{\alpha\in[4]}\vert\cF_{\alpha}\vert\leq 5\cdot n^{v(F)-1}\,.$$
    Due to the chosen hierarchy, this is a contradiction.
    \end{proof}

\end{document}
